\documentclass[10pt,aspectratio=169]{beamer}
\input{../../_en_preambulo_beamer}% Comandos y macros estadísticas compartidas para las presentaciones Beamer en inglés (Metropolis)

% Conjuntos numéricos básicos
\newcommand{\R}{\mathbb{R}}
\newcommand{\N}{\mathbb{N}}
\newcommand{\Q}{\mathbb{Q}}
\newcommand{\Z}{\mathbb{Z}}
\newcommand{\set}[1]{\left\{ #1 \right\}}
\newcommand{\abs}[1]{\left|#1\right|}
\newcommand{\norm}[1]{\left\| #1 \right\|}

% Comandos estadísticos y probabilísticos del curso
\newcommand{\Var}{\operatorname{Var}}
\newcommand{\cov}{\operatorname{Cov}}
\newcommand{\corr}{\rho}
\newcommand{\comb}[2]{\begin{pmatrix} #1 \\ #2 \end{pmatrix}}
\newcommand{\s}{\sigma}
\newcommand{\particion}{\mathfrak{P}}
\newcommand{\card}[1]{\#\left( #1 \right)}
\newcommand{\seq}[3]{\set{#1}_{#2}^{#3}}
\newcommand{\E}{\mathbb{E}}
\newcommand{\Prob}{\mathbb{P}}

% Entornos pedagógicos y cajas en inglés académico natural (soportando tanto nombres en inglés como en español)
\newenvironment{discussion}{%
  \begin{alertblock}{Discussion Question}%
}{%
  \end{alertblock}%
}
\newenvironment{discusion}{%
  \begin{alertblock}{Discussion Question}%
}{%
  \end{alertblock}%
}

\newenvironment{reflection}{%
  \begin{alertblock}{Classroom Reflection}%
}{%
  \end{alertblock}%
}
\newenvironment{reflexion}{%
  \begin{alertblock}{Classroom Reflection}%
}{%
  \end{alertblock}%
}

\newenvironment{paradox}{%
  \begin{alertblock}{Exploratory Paradox}%
}{%
  \end{alertblock}%
}
\newenvironment{paradoja}{%
  \begin{alertblock}{Exploratory Paradox}%
}{%
  \end{alertblock}%
}

\newenvironment{motivation}{%
  \begin{exampleblock}{Motivation: Data Science in Practice}%
}{%
  \end{exampleblock}%
}
\newenvironment{motivacion}{%
  \begin{exampleblock}{Motivation: Data Science in Practice}%
}{%
  \end{exampleblock}%
}

\newenvironment{quickchallenge}{%
  \begin{block}{Quick Team Challenge}%
}{%
  \end{block}%
}
\newenvironment{retoclase}{%
  \begin{block}{Quick Team Challenge}%
}{%
  \end{block}%
}

\title{Section 07.07: Proportion Tests}\subtitle{One and two binary populations}
\author[J. Castillo Colmenares]{Juliho Castillo Colmenares}\institute{Tecnológico de Monterrey}\date{}
\begin{document}
\begin{frame}[plain]\titlepage\end{frame}
\begin{frame}{Roadmap}\small\begin{enumerate}\item One proportion.\item Difference of two proportions.\item Conditions and decisions.\end{enumerate}\end{frame}
\begin{frame}{Source and scope}\small This section tests success/failure proportions for categorical or binary variables.\begin{block}{Guiding question}Is the observed proportion compatible with a reference value?\end{block}\end{frame}
\begin{frame}{Binomial model}\small Let $X$ count successes in $n$ trials with probability $p$. The sample proportion is $\hat p=X/n$.\end{frame}
\begin{frame}{One-proportion hypotheses}\small\[H_0:p=p_0\qquad\text{versus an alternative about }p.\]\end{frame}
\begin{frame}{One-proportion statistic}\small\[Z=\frac{\hat p-p_0}{\sqrt{p_0(1-p_0)/n}}.\]\end{frame}
\begin{frame}{Normal approximation}\small The approximation requires $np_0\geq5$ and $n(1-p_0)\geq5$. Otherwise consider an exact method.\end{frame}
\begin{frame}{Two-proportion test}\small For two independent populations, equality is $H_0:p_1-p_2=0$. Observe $\hat p_1$ and $\hat p_2$.\end{frame}
\begin{frame}{Pooled proportion}\small Under $H_0:p_1=p_2=p$, estimate\[\bar p=\frac{X_1+X_2}{n_1+n_2}.\]\end{frame}
\begin{frame}{Difference statistic}\small\[Z=\frac{\hat p_1-\hat p_2}{\sqrt{\bar p(1-\bar p)(1/n_1+1/n_2)}}.\]\end{frame}
\begin{frame}{Procedure I}\small\begin{enumerate}\item Define success and failure.\item State hypotheses.\item Check approximation conditions.\end{enumerate}\end{frame}
\begin{frame}{Procedure II}\small\begin{enumerate}\setcounter{enumi}{3}\item Compute proportions and statistic.\item Obtain critical value or $p$-value.\item Communicate the decision.\end{enumerate}\end{frame}
\begin{frame}{Alternative direction}\small Tails depend on whether the target is larger, smaller, or different. The sign of $\hat p-p_0$ guides interpretation.\end{frame}
\begin{frame}{Application: one proportion}\small A sample contains $X$ successes among $n$ observations. Test $H_0:p=p_0$ at level $\alpha$.\end{frame}
\begin{frame}{Application: calculation}\small Compute $\hat p=X/n$ and substitute into $Z$. The denominator is evaluated under $p_0$, not $\hat p$.\end{frame}
\begin{frame}{Application: decision}\small Compare $Z$ with the alternative's critical region or use the $p$-value.\end{frame}
\begin{frame}{Application: two populations}\small Two independent samples provide $X_1,n_1$ and $X_2,n_2$. Test equality using the pooled proportion.\end{frame}
\begin{frame}{Application: interpretation}\small Report a detectable difference together with proportions, sample sizes, and direction.\end{frame}
\begin{frame}{Common errors}\small\begin{itemize}\item Use $\hat p$ in the null standard error.\item Ignore approximation conditions.\item Confuse statistical and practical importance.\end{itemize}\end{frame}
\begin{frame}{Reporting template}\small\begin{block}{Report} ``$\hat p=\cdots$ (or $\hat p_1-\hat p_2=\cdots$), $Z=\cdots$, $p=\cdots$. At $\alpha=\cdots$, we [reject/do not reject] $H_0$.''\end{block}\end{frame}
\begin{frame}{Synthesis}\small\begin{itemize}\item One proportion uses the null value $p_0$.\item Two proportions use a pooled estimate under equality.\item Conditions support the normal approximation.\end{itemize}\end{frame}
\begin{frame}{Chapter connection}\small These tests prepare multiple-proportion comparisons and homogeneity tests.\end{frame}
\end{document}
